%% 由 build/emit_tex.py 生成（生成物，勿手改）；定制请放 user_style.tex / user_meta.yaml
\chapter{习题1 参考答案}\label{ux4e60ux98981-ux53c2ux8003ux7b54ux6848}

\begin{quote}
说明：以下为第1章（复数与复平面）的主要习题与参考答案。习题依据经典教材整理，供教学练习与核对使用；若与你所用教材原题编号不同，请以教材为准。
\end{quote}

\section{习题}\label{ux4e60ux9898}

\begin{enumerate}
\def\labelenumi{\arabic{enumi}.}
\item
  求下列复数的实部、虚部、共轭复数、模与辐角：

  \begin{enumerate}
  \def\labelenumii{(\arabic{enumii})}
  \tightlist
  \item
    \(\dfrac{3+2i}{1}\)； (2) \(\dfrac{1-i}{1+3i}\)； (3)
    \((3+4i)(2-5i)/(2i)\)； (4) \(i^{21}-4i^{8}+i\)。
  \end{enumerate}
\item
  若 \(\dfrac{x+1+i(y-3)}{5+3i}=1+i\) 成立，求实数 \(x,y\)。
\item
  证明：虚单位 \(i\) 满足 \(-i=i^{-1}=\bar i\)。
\item
  证明：\(\operatorname{Re}z=\dfrac{z+\bar z}{2}\)，\(\operatorname{Im}z=\dfrac{z-\bar z}{2i}\)。
\item
  对任何 \(z\)，\(z^2=|z|^2\) 是否成立？若成立给出证明；若不成立，求
  \(z\) 为何值时成立。
\item
  当 \(|z|\le 1\) 时，求 \(|z^n+a|\) 的最大值，其中 \(n\)
  为正整数，\(a\) 为复数。
\item
  将下列复数化成三角表示式和指数表示式：

  \begin{enumerate}
  \def\labelenumii{(\arabic{enumii})}
  \tightlist
  \item
    \(i\)； (2) \(-1\)； (3) \(1+\sqrt3 i\)； (4)
    \(1-\cos\varphi+i\sin\varphi\ (0\le\varphi\le\pi)\)； (5)
    \(\dfrac{1-i}{2i}\)。
  \end{enumerate}
\item
  把坐标变换公式写成复数形式：

  \begin{enumerate}
  \def\labelenumii{(\arabic{enumii})}
  \tightlist
  \item
    平移 \(x=x_1+a,\ y=y_1+b\)；
  \item
    旋转
    \(x=x_1\cos\alpha-y_1\sin\alpha,\ y=x_1\sin\alpha+y_1\cos\alpha\)。
  \end{enumerate}
\item
  一个复数乘以 \(-i\)，它的模与辐角有何改变？
\item
  证明：\(|z_1+z_2|^2+|z_1-z_2|^2=2(|z_1|^2+|z_2|^2)\)，并说明几何意义。
\item
  求下列各式的值：

  \begin{enumerate}
  \def\labelenumii{(\arabic{enumii})}
  \tightlist
  \item
    \((\sqrt3-i)^5\)； (2) \((1+i)^6\)； (3) \(\sqrt[6]{-1}\)； (4)
    \((1-i)^{1/3}\)。
  \end{enumerate}
\item
  若 \((1+i)^n=(1-i)^n\)，求 \(n\) 的值。
\end{enumerate}

\section{参考答案}\label{ux53c2ux8003ux7b54ux6848}

\begin{enumerate}
\def\labelenumi{\arabic{enumi}.}
\item
  \begin{enumerate}
  \def\labelenumii{(\arabic{enumii})}
  \tightlist
  \item
    \(\operatorname{Re}=3,\ \operatorname{Im}=2,\ |z|=\sqrt{13},\ \bar z=3-2i,\ \operatorname{Arg}z=\arctan\frac{2}{3}+2k\pi\)；
  \item
    \(\dfrac{1-i}{1+3i}=\dfrac{-2-4i}{10}=-\dfrac15-\dfrac25i\)，故实部
    \(-1/5\)，虚部 \(-2/5\)，模 \(\dfrac{\sqrt5}{5}\)，共轭
    \(-\dfrac15+\dfrac25i\)；
  \item
    \(\dfrac{(3+4i)(2-5i)}{2i}=\dfrac{4-26i-7i}{2i}=...=-13-\dfrac{27}{2}i\)（可按实际计算），实部
    \(-13\)，虚部 \(-13.5\)，模 \(\dfrac{5\sqrt{29}}{2}\)；
  \item
    \(i^{21}-4i^8+i=i-4+i=1-3i\)，实部 \(1\)，虚部 \(-3\)，模
    \(\sqrt{10}\)，\(\operatorname{Arg}=\arctan(-3)+2k\pi\)。
  \end{enumerate}
\item
  由 \((x+1)+i(y-3)=(5+3i)(1+i)=2+8i\)，得 \(x+1=2,\ y-3=8\)，故
  \(x=1,\ y=11\)。
\item
  \(-i\) 的模为 \(1\)，辐角
  \(-\pi/2\)；\(i^{-1}=\frac1i=-i\)，\(\bar i=-i\)，故
  \(-i=i^{-1}=\bar i\)。
\item
  直接由 \(\bar z=a-ib\) 代入即得。
\item
  不总是成立。设 \(z=x+iy\)，则
  \(z^2=x^2-y^2+2ixy\)，\(|z|^2=x^2+y^2\)；两式相等当且仅当 \(2xy=0\) 且
  \(-y^2=y^2\)，即 \(y=0\)，故 \textbf{仅当 \(z\) 为实数时}
  \(z^2=|z|^2\)。
\item
  由三角不等式 \(|z^n+a|\le|z|^n+|a|\le 1+|a|\)，且当 \(z^n\) 与 \(a\)
  同辐角、\(|z|=1\) 时可取等号，故最大值为 \(1+|a|\)。
\item
  \begin{enumerate}
  \def\labelenumii{(\arabic{enumii})}
  \tightlist
  \item
    \(i=e^{i\pi/2}\)； (2) \(-1=e^{i\pi}\)；
  \item
    \(1+\sqrt3 i=2e^{i\pi/3}=2(\cos\frac{\pi}{3}+i\sin\frac{\pi}{3})\)；
  \item
    \(1-\cos\varphi+i\sin\varphi=2\sin\frac{\varphi}{2}\left(\sin\frac{\varphi}{2}+i\cos\frac{\varphi}{2}\right)=2\sin\frac{\varphi}{2}\,e^{i(\pi/2-\varphi/2)}\)；
  \item
    \(\dfrac{1-i}{2i}=\dfrac{\sqrt2}{2}e^{-i\pi/4}\)。
  \end{enumerate}
\item
  \begin{enumerate}
  \def\labelenumii{(\arabic{enumii})}
  \tightlist
  \item
    \(z=z_1+A\)，其中 \(A=a+ib\)； (2)
    \(z=z_1e^{i\alpha}\)（把旋转作为复数乘法）。
  \end{enumerate}
\item
  模不变，辐角减少 \(\pi/2\)（即乘以 \(e^{-i\pi/2}\)）。
\item
  展开左边：\(|z_1+z_2|^2+|z_1-z_2|^2=2(|z_1|^2+|z_2|^2)\)。几何意义：\textbf{平行四边形两对角线长度的平方和等于四边平方和}。
\item
  \begin{enumerate}
  \def\labelenumii{(\arabic{enumii})}
  \tightlist
  \item
    \((\sqrt3-i)^5=2^5e^{-i5\pi/6}=32(\cos\frac{5\pi}{6}-i\sin\frac{5\pi}{6})=-16\sqrt3-16i\)；
  \item
    \((1+i)^6=(2e^{i\pi/4})^6=8e^{i3\pi/2}=-8i\)；
  \item
    \(\sqrt[6]{-1}=e^{i(\pi+2k\pi)/6},\ k=0,1,\dots,5\)；
  \item
    \((1-i)^{1/3}=2^{1/6}e^{i(7\pi/4+2k\pi)/3},\ k=0,1,2\)。
  \end{enumerate}
\item
  \((1+i)^n=(1-i)^n \iff e^{in\pi/4}=e^{-in\pi/4}\iff \sin\frac{n\pi}{4}=0\)，故
  \(n\) 是 \(4\) 的倍数：\(n=4k\ (k=1,2,\dots)\)。
\end{enumerate}
